\begin{lemma}[Any limit of the truncated Gaussian Riesz integrals is the corresponding coordinate of the operator]
Fix $n \geq 1$, let $\lambda$ be Lebesgue measure on $\mathbb{R}^n$, let
$\gamma_n$ be the Gaussian measure of \noderef{GaussMeasure}, and let
$K_{\gamma,j}$ be the kernel of \noderef{GaussRieszKernel}.  Let
\[
  T : L^1(\gamma_n;\mathbb{R}) \longrightarrow
      \mathcal{M}\bigl(\gamma_n;\mathbb{R}^n\bigr)
\]
be \emph{any} linear map into the $\gamma_n$-almost-everywhere equivalence
classes of measurable $\mathbb{R}^n$-valued functions such that
\begin{enumerate}
\item[(T1)] $T(\iota v) = R_\gamma v$ for every $v \in L^2(\gamma_n)$, with
$\iota$ the inclusion of \noderef{ToL1} and $R_\gamma$ the transform of
\noderef{RieszTransform}; and
\item[(T2)] $t\,\gamma_n\bigl(\{x : |Th(x)|_{\mathbb{R}^n} > t\}\bigr)
\leq 2\,\|h\|_{L^1(\gamma_n)}$ for every $h \in L^1(\gamma_n;\mathbb{R})$ and
every real $t > 0$.
\end{enumerate}
For a $\gamma_n$-integrable $h$, a coordinate $j$, a point $x$ and a real
$\varepsilon > 0$ write
$T_{j,\varepsilon}h(x) = \int_{|x-y|>\varepsilon}K_{\gamma,j}(x,y)h(y)\,dy$,
the $y$-integral against $\lambda$; for an element of $L^1(\gamma_n)$ this is
read through its canonical representative.

Let $f \in L^1(\gamma_n)$ and let $j \in \{1,\ldots,n\}$.  Then for
$\gamma_n$-almost every $x \in \mathbb{R}^n$ the following holds: for every
real number $L$,
\[
  \lim_{\varepsilon\downarrow0}T_{j,\varepsilon}f(x) = L
  \qquad\Longrightarrow\qquad
  L = (Tf)_j(x) ,
\]
the limit taken through $\varepsilon \in (0,\infty)$ as $\varepsilon \to 0$ and
$(Tf)_j(x)$ the $j$th coordinate of the canonical representative of $Tf$.

\emph{Identification only, and the implication form is the point.}  The
statement asserts nothing about whether the limit exists: it is vacuously true
at any $x$ where no $L$ is a limit.  That is deliberate.  Existence of the
limit is a separate assertion with separate inputs, and a statement combining
the two would be the conclusion that a covering node draws from this one
together with that one, so stating it here would leave that node nothing to do.
Keeping identification and existence apart is what makes each of them
independently citable.

\emph{Why quantifying over all $L$ loses nothing.}  Limits along
$\varepsilon \downarrow 0$ in $\mathbb{R}$ are unique, so at a given $x$ there
is at most one such $L$ and the universally quantified form says exactly "the
limit, if it exists, is $(Tf)_j(x)$".  The quantified form is used because it
needs no existence hypothesis and no partial function.

\emph{Why the conclusion is almost-everywhere.}  Two exceptional sets are
unavoidable, and both are controlled only in measure.  The first is where the
maximal truncation of the approximation error is large, controlled by
\noderef{GaussMaximalTruncationLocal}; the second is where $T$ of the error is
large, controlled by \noderef{CoordinateEndpointBound}.  Neither is empty for a
general $f \in L^1(\gamma_n)$, and the right-hand side $(Tf)_j$ is in any case
a coordinate of an almost-everywhere equivalence class, so no pointwise
statement about it is available.

\emph{Both (T1) and (T2) are assumed, because the inputs carry both.}  It is a
statement about every $T$ satisfying the two, and it is not vacuous because
\noderef{L1Extension} exhibits such a $T$.
\end{lemma}
\begin{proof}
\emph{Step 1: the approximants, and the error as a difference of classes.}
For each $k \in \mathbb{N}$ apply \noderef{GaussRieszTruncationSmoothSplit} to
$f$ and $k$.  It produces $\varphi_k$ and $g_k \in L^1(\gamma_n)$ with
$\varphi_k$ smooth, compactly supported and $\gamma_n$-integrable; with
\begin{equation*}
  \int_{\mathbb{R}^n}|g_k|\,d\gamma_n < \frac1{k+1} ; \tag{1}
\end{equation*}
with $f = \varphi_k + g_k$ $\gamma_n$-almost everywhere; and with, for every
$x$ and every real $\varepsilon > 0$, each of the three integrands Lebesgue
integrable on $\{y : |x-y| > \varepsilon\}$ and
\begin{equation*}
  T_{j,\varepsilon}f(x)
    = T_{j,\varepsilon}\varphi_k(x) + T_{j,\varepsilon}g_k(x) . \tag{2}
\end{equation*}
Write $\psi_k \in L^1(\gamma_n)$ for the class of $\varphi_k$, which exists
because $\varphi_k$ is $\gamma_n$-integrable.  The third clause says that the
canonical representative of $g_k$ agrees $\gamma_n$-almost everywhere with
$f - \varphi_k$, hence with the canonical representative of $f - \psi_k$; and
two elements of $L^1(\gamma_n)$ whose representatives agree almost everywhere
are equal, so
\begin{equation*}
  g_k = f - \psi_k \quad\text{in } L^1(\gamma_n) . \tag{3}
\end{equation*}
Also $\|g_k\|_{L^1(\gamma_n)} = \int|g_k|\,d\gamma_n < 1/(k+1)$ by $(1)$.

\emph{Step 2: two almost-everywhere facts, for all $k$ at once.}
For each $k$, \noderef{SmoothCoreTruncationBridge} applied to this $T$ with its
property (T1), to $\varphi_k$ --- smooth, compactly supported,
$\gamma_n$-integrable --- and to the coordinate $j$, gives a $\gamma_n$-null
set off which
\begin{equation*}
  \lim_{\varepsilon\downarrow0}T_{j,\varepsilon}\varphi_k(x)
    = (T\psi_k)_j(x) . \tag{4}
\end{equation*}
And $T$ is linear, so by $(3)$, $Tg_k = Tf - T\psi_k$ in
$\mathcal{M}(\gamma_n;\mathbb{R}^n)$; the canonical representative of a
difference of classes agrees $\gamma_n$-almost everywhere with the difference
of the representatives, so there is a $\gamma_n$-null set off which
\begin{equation*}
  (Tg_k)_j(x) = (Tf)_j(x) - (T\psi_k)_j(x) . \tag{5}
\end{equation*}
A countable union of $\gamma_n$-null sets is $\gamma_n$-null, so there is a
single $\gamma_n$-null set $N$ off which $(4)$ and $(5)$ hold \emph{for every}
$k \in \mathbb{N}$.  It therefore suffices to show that the set
\[
  B = \bigl\{x : \text{some real } L \text{ is a limit of }
     T_{j,\varepsilon}f(x) \text{ and } L \neq (Tf)_j(x)\bigr\}
\]
satisfies $\gamma_n(B \setminus N) = 0$, since then
$\gamma_n(B) \leq \gamma_n(B\setminus N) + \gamma_n(N) = 0$ by subadditivity.

\emph{Step 3: a two-way split in measure.}
Note first that the split is two-way and not three-way: $(4)$ identifies the
limit of the smooth part \emph{as} $(T\psi_k)_j(x)$ directly, so no separate
set on which that identification fails has to be introduced.

Let $x \in B \setminus N$, with witness $L$.  Then
$|L - (Tf)_j(x)| > 0$, so there is $m \in \mathbb{N}$ with
\begin{equation*}
  2\cdot\frac1{m+1} < \bigl|L - (Tf)_j(x)\bigr| , \tag{6}
\end{equation*}
and there is $r \in \mathbb{N}$ with $|x| \leq r+1$, for instance
$r = \lceil|x|\rceil$.  Fix any $k$.  By $(5)$ and the triangle inequality,
\[
  \bigl|L - (Tf)_j(x)\bigr|
    \leq \bigl|L - (T\psi_k)_j(x)\bigr| + \bigl|(Tg_k)_j(x)\bigr| ,
\]
so by $(6)$ at least one of the two terms on the right exceeds $1/(m+1)$.

Suppose $|(Tg_k)_j(x)| > 1/(m+1)$.  Then $x$ lies in
$F_{m,k} = \{x : |(Tg_k)_j(x)| > 1/(m+1)\}$.

Suppose instead $|L - (T\psi_k)_j(x)| > 1/(m+1)$.  By $(2)$ at each
$\varepsilon > 0$, $T_{j,\varepsilon}g_k(x)
 = T_{j,\varepsilon}f(x) - T_{j,\varepsilon}\varphi_k(x)$, and both terms on
the right converge as $\varepsilon \downarrow 0$ --- the first to $L$ by the
choice of $L$, the second to $(T\psi_k)_j(x)$ by $(4)$, available because
$x \notin N$.  Hence $T_{j,\varepsilon}g_k(x) \to L - (T\psi_k)_j(x)$ and so
$|T_{j,\varepsilon}g_k(x)| \to |L - (T\psi_k)_j(x)| > 1/(m+1)$.  A convergent
family whose limit exceeds $1/(m+1)$ eventually exceeds it, and every
neighbourhood of $0$ meets $(0,\infty)$, so there is a real $\varepsilon > 0$
with $|T_{j,\varepsilon}g_k(x)| > 1/(m+1)$.  Since a single term is at most the
supremum and $\mathrm{ofReal}$ is strictly increasing on $[0,\infty)$,
\[
  \mathrm{ofReal}\Bigl(\frac1{m+1}\Bigr)
    < \sup_{\varepsilon>0}\mathrm{ofReal}
      \bigl|T_{j,\varepsilon}g_k(x)\bigr| ,
\]
the supremum in $[0,+\infty]$; together with $|x| \leq r+1$ this puts $x$ in
\[
  E_{r,m,k} = \Bigl\{x : |x| \leq r+1 \text{ and }
    \mathrm{ofReal}\bigl(\tfrac1{m+1}\bigr)
      < \sup_{\varepsilon>0}\mathrm{ofReal}\bigl|T_{j,\varepsilon}g_k(x)\bigr|
  \Bigr\} .
\]
As $k$ was arbitrary,
$x \in \bigcap_{k}\bigl(E_{r,m,k} \cup F_{m,k}\bigr)$, and therefore
\[
  B \setminus N \;\subseteq\;
    \bigcup_{m\in\mathbb{N}}\bigcup_{r\in\mathbb{N}}
      \bigcap_{k\in\mathbb{N}}\bigl(E_{r,m,k} \cup F_{m,k}\bigr) .
\]
By countable subadditivity it suffices to show
$\gamma_n\bigl(\bigcap_k(E_{r,m,k} \cup F_{m,k})\bigr) = 0$ for each fixed $r$
and $m$.

Fix $r$ and $m$ and put $t_m = \mathrm{ofReal}(1/(m+1)) \neq 0$.  Apply
\noderef{GaussMaximalTruncationLocal} to this $T$ with its two properties, to
the coordinate $j$ and to the radius $R := r+1 \geq 1$; it supplies
$C_r \in (0,\infty)$, produced before its input is quantified, with
$\gamma_n(E_{r,m,k}) \leq C_r\cdot\mathrm{ofReal}\bigl(\int|g_k|d\gamma_n\bigr)/t_m$
for every $k$.  Apply \noderef{CoordinateEndpointBound} to this $T$ with its two
properties, to the input $g_k$, the coordinate $j$ and the real level
$1/(m+1) > 0$; it gives
$t_m\cdot\gamma_n(F_{m,k}) \leq \mathrm{ofReal}(2\|g_k\|_{L^1(\gamma_n)})$,
hence $\gamma_n(F_{m,k}) \leq
 \mathrm{ofReal}(2\|g_k\|_{L^1(\gamma_n)})/t_m$ because $t_m$ is neither $0$
nor $+\infty$.  Let
$D = C_r/t_m + \mathrm{ofReal}(2)/t_m$, which is not $+\infty$ because $C_r$ is
not and $t_m$ is not $0$.  Using $ab/c = (a/c)b$ in $[0,+\infty]$, then $(1)$
and Step 1's norm identity, both bounds become at most
$D\cdot\mathrm{ofReal}(1/(k+1))$.  So for every $k$, by monotonicity and
subadditivity,
\[
  \gamma_n\Bigl(\bigcap_{i}(E_{r,m,i} \cup F_{m,i})\Bigr)
    \leq \gamma_n(E_{r,m,k}) + \gamma_n(F_{m,k})
    \leq (D + D)\cdot\mathrm{ofReal}\Bigl(\frac1{k+1}\Bigr) .
\]
Now $1/(k+1) \to 0$, so $\mathrm{ofReal}(1/(k+1)) \to 0$ in $[0,+\infty]$, and
since $D + D \neq +\infty$ also $(D+D)\cdot\mathrm{ofReal}(1/(k+1)) \to 0$.  A
quantity in $[0,+\infty]$ bounded above by every term of a sequence converging
to $0$ is at most $0$, so the intersection is $\gamma_n$-null, completing the
proof.
\end{proof}
